%%%%%%%%%%%%%%%%%%%%%%%%%%%%
%%%%%     PREAMBLE     %%%%%
%%%%%%%%%%%%%%%%%%%%%%%%%%%%

\documentclass[11pt,twocolumn]{article}

%Define Margins using Geometry Package
\usepackage[top=1.0in,
bottom=1.0in,
right=1.0in,
left = 1.0in]{geometry}


%math
\usepackage{mathtools}


%Set up links and internal references
\usepackage[colorlinks=true,linkcolor=blue, urlcolor=blue]{hyperref} %hide the ugly red boxes
\usepackage[noabbrev]{cleveref} %don't abbreviate ``figure'', etc.
\usepackage{url} %for urls

%Remove natural indentation for itemize and enumerate environments
\usepackage{enumitem} %package to control itemize/enumerate behavior
\setlist[itemize]{leftmargin=*} %for itemize
\setlist[enumerate]{leftmargin=*} %for enumerate
%Swap out the itemize symbol for an N-dash rather than a bullet (does not require enumitem package)
%\def\labelitemi{--}

\usepackage{multicol} %Allow for two columns in the middle of a paper. 


%Set up fancy header and footer
\usepackage{fancyhdr}
\pagestyle{fancy}
%Fancy Header Content
\fancyhead[L]{ME EN 497R}
\fancyhead[C]{}
\fancyhead[R]{Rotor Design Theory Exploration: Part I}
%Fancy Footer Content
\fancyfoot[L,C]{} %make bottom left and center empty
\fancyfoot[R]{\thepage}


% Customize formatting for section headers
\renewcommand{\thesection}{Problem \arabic{section}}
\renewcommand{\thesubsection}{\arabic{section}.\alph{subsection}}

%%%%%%%%%%%%%%%%%%%%%%%%%%%%%%%%
%%%%%     END PREAMBLE     %%%%%
%%%%%%%%%%%%%%%%%%%%%%%%%%%%%%%%


%BEGIN Document Environment
\begin{document}



%%%%%%%%%%%%%%%%%%%%%%%%%%%%%%%%%%%%%%%%%%%%%%%%%%%%%%%%%%%%%%%%%%%%%%
%                                                                    %
%                               Reading                              %
%                                                                    %
%%%%%%%%%%%%%%%%%%%%%%%%%%%%%%%%%%%%%%%%%%%%%%%%%%%%%%%%%%%%%%%%%%%%%%
\section*{Relevant Reading}
\label{sec:reading}

\begin{itemize}
	\item \href{http://flowlab.groups.et.byu.net/me415/flight.pdf}{Flight Vehicle Design Book}
	\begin{itemize}
		\item Chapter 2: Fundamentals (13 pages)
		\item Chapter 6.3: Propellers (4 pages)
	\end{itemize}

\end{itemize}




%%%%%%%%%%%%%%%%%%%%%%%%%%%%%%%%%%%%%%%%%%%%%%%%%%%%%%%%%%%%%%%%%%%%%%
%                                                                    %
%                             Vocabulary                             %
%                                                                    %
%%%%%%%%%%%%%%%%%%%%%%%%%%%%%%%%%%%%%%%%%%%%%%%%%%%%%%%%%%%%%%%%%%%%%%

\section{Vocabulary}
\label{sec:vocab}

Explain the following terms; making sure to use sufficient detail, including any math or helpful figures.
In some cases, these terms are simple one sentence definitions, in others, you should include several paragraphs to explain them fully.

\paragraph{Airfoil Geometry Terms}
\begin{itemize}
	\item Chord
	\item Leading Edge
	\item Trailing Edge
\end{itemize}

\paragraph{Rotor Geometry}
\begin{itemize}
	\item Hub Radius
	\item Tip Radius
	\item Twist
	\item Pitch
	\item Precone
	\item Sweep
\end{itemize}

\paragraph{Non-dimensional Numbers}
\begin{itemize}
	\item Reynolds Number
	\item Mach Number
	\item Advance ratio
	\item Tip Speed Ratio
	\item Solidity
\end{itemize}


\paragraph{Forces and Moments}
\begin{itemize}
	\item Lift
	\item Drag
	\item Pitching Moment
	\item Lift and Drag Polars
	\begin{itemize}
		\item Angle of Attack
		\item Zero Lift angle of attack
		\item Lift Curve Slope
		\item Stall
	\end{itemize}
\end{itemize}




%%%%%%%%%%%%%%%%%%%%%%%%%%%%%%%%%%%%%%%%%%%%%%%%%%%%%%%%%%%%%%%%%%%%%%
%                                                                    %
%                               Explore                              %
%                                                                    %
%%%%%%%%%%%%%%%%%%%%%%%%%%%%%%%%%%%%%%%%%%%%%%%%%%%%%%%%%%%%%%%%%%%%%%
\section{Exploration}
\label{sec:exploration}

Complete the following exploration.

\subsection{Prerequisites}
\label{ssec:prereqs}

\begin{enumerate}[label=\roman*.]
	\item \href{https://flow.byu.edu/CCBlade.jl/stable/}{Install CCBlade.jl}
	\item Complete the \href{https://flow.byu.edu/CCBlade.jl/stable/tutorial/}{Quick Start Tutorial}
\end{enumerate}


\subsection{Load Distributions and Rotor Efficiency}

Blade Element Momentum Theory (BEMT) allows for the analysis of propellers. wind turbines, and helicopter blades. 
Choose one that interests you and implement it in CCBlade.jl.

\begin{enumerate}[label=\roman*.]
	\item Plot the section loads (normal and tangential loads) vs radius. Analyze and explain the shape of the load distribution. Why does the loading increase? Why are there decreases? 
	\item Compare the thrust, torque, and root bending moment vs varying rotor tip radius.
	\item Compare the thrust, torque, and root bending moment vs varying the rotor section chord length.
\end{enumerate}





\end{document}